% Methods fragment for peer-reviewed papers using FLBacktest
% Include with: % Methods fragment for peer-reviewed papers using FLBacktest
% Include with: % Methods fragment for peer-reviewed papers using FLBacktest
% Include with: % Methods fragment for peer-reviewed papers using FLBacktest
% Include with: \input{path/to/methods-backtest.tex}
% Or copy subsections into the host paper's Methods section.

\section{Methods}
\label{sec:methods-backtest}

\subsection{Scope}

The analysis evaluates harvest control rules (HCRs) with a historical
closed-loop backtest implemented in the FLR package \texttt{FLBacktest}.
The procedure separates two questions that are often conflated: whether the
assessment and advice rule are robust under historical uncertainty, and
whether realised management outcomes differ because advice was not followed.

\subsection{Overview of the backtest}

For each operating-model stock and each advice year \(t\) in a historical
window \([t_0,T]\):
\begin{enumerate}
  \item \textbf{Perceive.} Form a management-procedure view of stock status
        using only information available at \(t\) (perfect information,
        short-cut assessment error, or a full re-assessment / peel).
  \item \textbf{Advise.} Apply the HCR in force at \(t\), with that year's
        reference points, to obtain catch (or \(F\)) advice for year
        \(t{+}1\).
  \item \textbf{Implement.} Project the operating model forward under the
        advised control (optionally with implementation error).
  \item \textbf{Compare.} Contrast simulated SSB, \(F\), and catch with the
        realised historical path over the same years.
\end{enumerate}
Repeating these steps for successive \(t\) builds a counterfactual path under
consistent adherence to advice as it was understood at the time
\citep{Punt2016MSE,Butterworth2014}.

\subsection{Operating model}

The operating model (OM) is an age-structured \texttt{FLStock} conditioned on
a historical assessment (or a set of plausible assessments). A stock--recruitment
relationship is represented as an \texttt{FLBRP} equilibrium for short-term
forecasts and closed-loop projections \citep{Kell2007FLR}. Multiple OMs may be
retained for robustness (e.g.\ alternative natural mortality or productivity).
Stochastic recruitment uses multiplicative deviations; papers should state
whether deviations are aleatory or include epistemic trend components.

\subsection{Observation and estimation error}

Three nested observation models are available.

\paragraph{Perfect information.}
The management procedure observes true OM status without error. This isolates
HCR behaviour from assessment error and is a diagnostic, not a realistic advice
process.

\paragraph{Short-cut assessment error.}
Multiplicative error is imposed on the HCR metric (typically SSB), conditioned
on retrospective peels. An ARIMA (or AR(1)) model fitted to retrospective bias
is simulated across years and iterations
\citep{ICES2013EqSim,Kell2016}, avoiding a full re-assessment each year while
retaining empirically motivated error structure.

\paragraph{Full feedback.}
Each year the management procedure updates an assessment from contemporaneous
catch and survey data generated from the OM, then applies the HCR to the
estimates, jointly including observation, estimation, and advice error
\citep{Punt2016MSE}.

\subsection{Harvest control rule}

The base HCR is the ICES Category~1 MSY hockey-stick: \(F=F_{\mathrm{MSY}}\)
above \(MSYB_{\mathrm{trigger}}\), reduced linearly towards \(B_{\lim}\), and
zero at or below \(B_{\lim}\) \citep{ICES2021Advice}. Catch advice is obtained
by short-term projection at the advised \(F\). Optional year-to-year TAC
stability constraints may apply when SSB is above a stated threshold.
Reference points are those used for advice in year \(t\); benchmark years are
structural breaks when reference points are redefined.

\subsection{Closed-loop projection}

Given advised catch (or \(F\)) for year \(t{+}1\), the OM is projected with
FLR forward methods under the chosen recruitment deviations. Management and
data lags are stated explicitly. Perfect-management counterfactuals set
realised catch equal to advised catch; implementation error may be added when
required.

\subsection{Performance metrics}

Realised and counterfactual histories are compared using:
\begin{itemize}
  \item \textbf{Biological status} --- time (or probability) below
        \(B_{\lim}\) and \(MSYB_{\mathrm{trigger}}\);
  \item \textbf{Exploitation} --- \(F\) relative to \(F_{\mathrm{MSY}}\);
  \item \textbf{Yield and stability} --- mean/cumulative catch and interannual
        variability;
  \item \textbf{Risk trade-offs} --- yield versus probability of falling below
        limit reference points \citep{Punt2016MSE}.
\end{itemize}

\subsection{Software}

Analyses use FLR (\texttt{FLCore}, \texttt{FLBRP}, \texttt{FLasher},
\texttt{ggplotFL}) and \texttt{FLBacktest} (\texttt{hcrICES}/\texttt{backtest},
\texttt{shortcut}, \texttt{retroBias}, \texttt{backtestResults}). Conditioned
OMs, scenario grids, seeds, and session information are archived so that
tables and figures can be regenerated from saved results.

% Or copy subsections into the host paper's Methods section.

\section{Methods}
\label{sec:methods-backtest}

\subsection{Scope}

The analysis evaluates harvest control rules (HCRs) with a historical
closed-loop backtest implemented in the FLR package \texttt{FLBacktest}.
The procedure separates two questions that are often conflated: whether the
assessment and advice rule are robust under historical uncertainty, and
whether realised management outcomes differ because advice was not followed.

\subsection{Overview of the backtest}

For each operating-model stock and each advice year \(t\) in a historical
window \([t_0,T]\):
\begin{enumerate}
  \item \textbf{Perceive.} Form a management-procedure view of stock status
        using only information available at \(t\) (perfect information,
        short-cut assessment error, or a full re-assessment / peel).
  \item \textbf{Advise.} Apply the HCR in force at \(t\), with that year's
        reference points, to obtain catch (or \(F\)) advice for year
        \(t{+}1\).
  \item \textbf{Implement.} Project the operating model forward under the
        advised control (optionally with implementation error).
  \item \textbf{Compare.} Contrast simulated SSB, \(F\), and catch with the
        realised historical path over the same years.
\end{enumerate}
Repeating these steps for successive \(t\) builds a counterfactual path under
consistent adherence to advice as it was understood at the time
\citep{Punt2016MSE,Butterworth2014}.

\subsection{Operating model}

The operating model (OM) is an age-structured \texttt{FLStock} conditioned on
a historical assessment (or a set of plausible assessments). A stock--recruitment
relationship is represented as an \texttt{FLBRP} equilibrium for short-term
forecasts and closed-loop projections \citep{Kell2007FLR}. Multiple OMs may be
retained for robustness (e.g.\ alternative natural mortality or productivity).
Stochastic recruitment uses multiplicative deviations; papers should state
whether deviations are aleatory or include epistemic trend components.

\subsection{Observation and estimation error}

Three nested observation models are available.

\paragraph{Perfect information.}
The management procedure observes true OM status without error. This isolates
HCR behaviour from assessment error and is a diagnostic, not a realistic advice
process.

\paragraph{Short-cut assessment error.}
Multiplicative error is imposed on the HCR metric (typically SSB), conditioned
on retrospective peels. An ARIMA (or AR(1)) model fitted to retrospective bias
is simulated across years and iterations
\citep{ICES2013EqSim,Kell2016}, avoiding a full re-assessment each year while
retaining empirically motivated error structure.

\paragraph{Full feedback.}
Each year the management procedure updates an assessment from contemporaneous
catch and survey data generated from the OM, then applies the HCR to the
estimates, jointly including observation, estimation, and advice error
\citep{Punt2016MSE}.

\subsection{Harvest control rule}

The base HCR is the ICES Category~1 MSY hockey-stick: \(F=F_{\mathrm{MSY}}\)
above \(MSYB_{\mathrm{trigger}}\), reduced linearly towards \(B_{\lim}\), and
zero at or below \(B_{\lim}\) \citep{ICES2021Advice}. Catch advice is obtained
by short-term projection at the advised \(F\). Optional year-to-year TAC
stability constraints may apply when SSB is above a stated threshold.
Reference points are those used for advice in year \(t\); benchmark years are
structural breaks when reference points are redefined.

\subsection{Closed-loop projection}

Given advised catch (or \(F\)) for year \(t{+}1\), the OM is projected with
FLR forward methods under the chosen recruitment deviations. Management and
data lags are stated explicitly. Perfect-management counterfactuals set
realised catch equal to advised catch; implementation error may be added when
required.

\subsection{Performance metrics}

Realised and counterfactual histories are compared using:
\begin{itemize}
  \item \textbf{Biological status} --- time (or probability) below
        \(B_{\lim}\) and \(MSYB_{\mathrm{trigger}}\);
  \item \textbf{Exploitation} --- \(F\) relative to \(F_{\mathrm{MSY}}\);
  \item \textbf{Yield and stability} --- mean/cumulative catch and interannual
        variability;
  \item \textbf{Risk trade-offs} --- yield versus probability of falling below
        limit reference points \citep{Punt2016MSE}.
\end{itemize}

\subsection{Software}

Analyses use FLR (\texttt{FLCore}, \texttt{FLBRP}, \texttt{FLasher},
\texttt{ggplotFL}) and \texttt{FLBacktest} (\texttt{hcrICES}/\texttt{backtest},
\texttt{shortcut}, \texttt{retroBias}, \texttt{backtestResults}). Conditioned
OMs, scenario grids, seeds, and session information are archived so that
tables and figures can be regenerated from saved results.

% Or copy subsections into the host paper's Methods section.

\section{Methods}
\label{sec:methods-backtest}

\subsection{Scope}

The analysis evaluates harvest control rules (HCRs) with a historical
closed-loop backtest implemented in the FLR package \texttt{FLBacktest}.
The procedure separates two questions that are often conflated: whether the
assessment and advice rule are robust under historical uncertainty, and
whether realised management outcomes differ because advice was not followed.

\subsection{Overview of the backtest}

For each operating-model stock and each advice year \(t\) in a historical
window \([t_0,T]\):
\begin{enumerate}
  \item \textbf{Perceive.} Form a management-procedure view of stock status
        using only information available at \(t\) (perfect information,
        short-cut assessment error, or a full re-assessment / peel).
  \item \textbf{Advise.} Apply the HCR in force at \(t\), with that year's
        reference points, to obtain catch (or \(F\)) advice for year
        \(t{+}1\).
  \item \textbf{Implement.} Project the operating model forward under the
        advised control (optionally with implementation error).
  \item \textbf{Compare.} Contrast simulated SSB, \(F\), and catch with the
        realised historical path over the same years.
\end{enumerate}
Repeating these steps for successive \(t\) builds a counterfactual path under
consistent adherence to advice as it was understood at the time
\citep{Punt2016MSE,Butterworth2014}.

\subsection{Operating model}

The operating model (OM) is an age-structured \texttt{FLStock} conditioned on
a historical assessment (or a set of plausible assessments). A stock--recruitment
relationship is represented as an \texttt{FLBRP} equilibrium for short-term
forecasts and closed-loop projections \citep{Kell2007FLR}. Multiple OMs may be
retained for robustness (e.g.\ alternative natural mortality or productivity).
Stochastic recruitment uses multiplicative deviations; papers should state
whether deviations are aleatory or include epistemic trend components.

\subsection{Observation and estimation error}

Three nested observation models are available.

\paragraph{Perfect information.}
The management procedure observes true OM status without error. This isolates
HCR behaviour from assessment error and is a diagnostic, not a realistic advice
process.

\paragraph{Short-cut assessment error.}
Multiplicative error is imposed on the HCR metric (typically SSB), conditioned
on retrospective peels. An ARIMA (or AR(1)) model fitted to retrospective bias
is simulated across years and iterations
\citep{ICES2013EqSim,Kell2016}, avoiding a full re-assessment each year while
retaining empirically motivated error structure.

\paragraph{Full feedback.}
Each year the management procedure updates an assessment from contemporaneous
catch and survey data generated from the OM, then applies the HCR to the
estimates, jointly including observation, estimation, and advice error
\citep{Punt2016MSE}.

\subsection{Harvest control rule}

The base HCR is the ICES Category~1 MSY hockey-stick: \(F=F_{\mathrm{MSY}}\)
above \(MSYB_{\mathrm{trigger}}\), reduced linearly towards \(B_{\lim}\), and
zero at or below \(B_{\lim}\) \citep{ICES2021Advice}. Catch advice is obtained
by short-term projection at the advised \(F\). Optional year-to-year TAC
stability constraints may apply when SSB is above a stated threshold.
Reference points are those used for advice in year \(t\); benchmark years are
structural breaks when reference points are redefined.

\subsection{Closed-loop projection}

Given advised catch (or \(F\)) for year \(t{+}1\), the OM is projected with
FLR forward methods under the chosen recruitment deviations. Management and
data lags are stated explicitly. Perfect-management counterfactuals set
realised catch equal to advised catch; implementation error may be added when
required.

\subsection{Performance metrics}

Realised and counterfactual histories are compared using:
\begin{itemize}
  \item \textbf{Biological status} --- time (or probability) below
        \(B_{\lim}\) and \(MSYB_{\mathrm{trigger}}\);
  \item \textbf{Exploitation} --- \(F\) relative to \(F_{\mathrm{MSY}}\);
  \item \textbf{Yield and stability} --- mean/cumulative catch and interannual
        variability;
  \item \textbf{Risk trade-offs} --- yield versus probability of falling below
        limit reference points \citep{Punt2016MSE}.
\end{itemize}

\subsection{Software}

Analyses use FLR (\texttt{FLCore}, \texttt{FLBRP}, \texttt{FLasher},
\texttt{ggplotFL}) and \texttt{FLBacktest} (\texttt{hcrICES}/\texttt{backtest},
\texttt{shortcut}, \texttt{retroBias}, \texttt{backtestResults}). Conditioned
OMs, scenario grids, seeds, and session information are archived so that
tables and figures can be regenerated from saved results.

% Or copy subsections into the host paper's Methods section.

\section{Methods}
\label{sec:methods-backtest}

\subsection{Scope}

The analysis evaluates harvest control rules (HCRs) with a historical
closed-loop backtest implemented in the FLR package \texttt{FLBacktest}.
The procedure separates two questions that are often conflated: whether the
assessment and advice rule are robust under historical uncertainty, and
whether realised management outcomes differ because advice was not followed.

\subsection{Overview of the backtest}

For each operating-model stock and each advice year \(t\) in a historical
window \([t_0,T]\):
\begin{enumerate}
  \item \textbf{Perceive.} Form a management-procedure view of stock status
        using only information available at \(t\) (perfect information,
        short-cut assessment error, or a full re-assessment / peel).
  \item \textbf{Advise.} Apply the HCR in force at \(t\), with that year's
        reference points, to obtain catch (or \(F\)) advice for year
        \(t{+}1\).
  \item \textbf{Implement.} Project the operating model forward under the
        advised control (optionally with implementation error).
  \item \textbf{Compare.} Contrast simulated SSB, \(F\), and catch with the
        realised historical path over the same years.
\end{enumerate}
Repeating these steps for successive \(t\) builds a counterfactual path under
consistent adherence to advice as it was understood at the time
\citep{Punt2016MSE,Butterworth2014}.

\subsection{Operating model}

The operating model (OM) is an age-structured \texttt{FLStock} conditioned on
a historical assessment (or a set of plausible assessments). A stock--recruitment
relationship is represented as an \texttt{FLBRP} equilibrium for short-term
forecasts and closed-loop projections \citep{Kell2007FLR}. Multiple OMs may be
retained for robustness (e.g.\ alternative natural mortality or productivity).
Stochastic recruitment uses multiplicative deviations; papers should state
whether deviations are aleatory or include epistemic trend components.

\subsection{Observation and estimation error}

Three nested observation models are available.

\paragraph{Perfect information.}
The management procedure observes true OM status without error. This isolates
HCR behaviour from assessment error and is a diagnostic, not a realistic advice
process.

\paragraph{Short-cut assessment error.}
Multiplicative error is imposed on the HCR metric (typically SSB), conditioned
on retrospective peels. An ARIMA (or AR(1)) model fitted to retrospective bias
is simulated across years and iterations
\citep{ICES2013EqSim,Kell2016}, avoiding a full re-assessment each year while
retaining empirically motivated error structure.

\paragraph{Full feedback.}
Each year the management procedure updates an assessment from contemporaneous
catch and survey data generated from the OM, then applies the HCR to the
estimates, jointly including observation, estimation, and advice error
\citep{Punt2016MSE}.

\subsection{Harvest control rule}

The base HCR is the ICES Category~1 MSY hockey-stick: \(F=F_{\mathrm{MSY}}\)
above \(MSYB_{\mathrm{trigger}}\), reduced linearly towards \(B_{\lim}\), and
zero at or below \(B_{\lim}\) \citep{ICES2021Advice}. Catch advice is obtained
by short-term projection at the advised \(F\). Optional year-to-year TAC
stability constraints may apply when SSB is above a stated threshold.
Reference points are those used for advice in year \(t\); benchmark years are
structural breaks when reference points are redefined.

\subsection{Closed-loop projection}

Given advised catch (or \(F\)) for year \(t{+}1\), the OM is projected with
FLR forward methods under the chosen recruitment deviations. Management and
data lags are stated explicitly. Perfect-management counterfactuals set
realised catch equal to advised catch; implementation error may be added when
required.

\subsection{Performance metrics}

Realised and counterfactual histories are compared using:
\begin{itemize}
  \item \textbf{Biological status} --- time (or probability) below
        \(B_{\lim}\) and \(MSYB_{\mathrm{trigger}}\);
  \item \textbf{Exploitation} --- \(F\) relative to \(F_{\mathrm{MSY}}\);
  \item \textbf{Yield and stability} --- mean/cumulative catch and interannual
        variability;
  \item \textbf{Risk trade-offs} --- yield versus probability of falling below
        limit reference points \citep{Punt2016MSE}.
\end{itemize}

\subsection{Software}

Analyses use FLR (\texttt{FLCore}, \texttt{FLBRP}, \texttt{FLasher},
\texttt{ggplotFL}) and \texttt{FLBacktest} (\texttt{hcrICES}/\texttt{backtest},
\texttt{shortcut}, \texttt{retroBias}, \texttt{backtestResults}). Conditioned
OMs, scenario grids, seeds, and session information are archived so that
tables and figures can be regenerated from saved results.
